\section{Materiales y Metodología}

\subsection{Materiales}

\begin{itemize}
	\item Sistema embebido: Arduino, Raspberry Pi Pico, ESP32 o STM32.
	\item Protoboard
	\item Cables de conexión
	\item Resistencias de distintos valores
	\item Pantalla LCD 16x2 con I2C
	\item Sensor de $CO_2$
	\item Sensor de temperatura ambiente 
	\item Sensor de humedad ambiente
\end{itemize}

\subsection{Metodología}

El estudiante junto con su grupo de trabajo colaborativo va a seleccionar un sistema embebido (Arduino, Raspberry Pi Pico, ESP32 o STM32) que permita dar solución al problema: \textit{Integración de Robótica Autónoma en Sistemas de Automatización para Procesos Ambientales}.

Se debe desarrollar la siguiente implementación:
\begin{enumerate}
	\item Conecte y configure una Pantalla LCD 16x2 con módulo conversor de paralelo a I2C, presentando el mensaje ``Sistema Robótico Automatizado para Procesos Ambientales, sensores'' realizando un scroll del texto de derecha a izquierda. 
	\item Configure el sensor de $CO_2$ de tal forma que se presente el valor de la medición en la pantalla
	LCD 16x2 I2C con su respectiva unidad de medida; integre en el mismo código a los sensores de 
	temperatura y humedad, teniendo en cuenta las configuraciones correspondientes para cada uno de estos. 
	Las tres magnitudes físicas se deben ver en la misma pantalla, puede ser de forma simultánea o con un 
	retardo de tiempo entre estos, realice la toma de diez cambios significativos para las magnitudes 
	físicas en una tabla con su respectiva evidencia, para el $CO_2$ puede apoyarse del gas de un briket.
\end{enumerate}
